\section{Agent-Based Model Results: A Synthetic-Population Stress Test}\label{sec:abm}

This section is a stress test of the household-choice hypothesis, not the paper's main estimate. And the first thing to report about it is what its own cross-design leg establishes about its scope, because that result governs how everything below should be read. Re-running the same decision rule on real Freddie Mac structural covariates recovers 59.3\% of the benchmark on the committed draw and 60.2\% averaged over fifty seeds---above the 50\% threshold fixed before that run at which this section's paradigm reading would be undercut. It is undercut on all fifty of them. That cross-design leg runs under the recalibration rule the model already applies elsewhere. Recovery rises to 76.3\% once the sample is reweighted to the SOMA book's own coupon composition. Three things bound that reading and travel with it wherever it appears. First, the same reweight moves the frozen-calibration leg the other way, down to 12.6\%. Second, the covariates swapped in are structural only, since the loan records contain no analogue for mobility desire or loss aversion. So no leg here replaces a synthetic behavioral population with a real one. Third, the verdict the run supports is that the contrast cannot, on that test alone, be cleanly attributed to modeling paradigm rather than to calibration and data source (Section~\ref{sec:robustness-crossdesign}). What this section falsifies is therefore a \emph{synthetic-population} household-choice specification: no specification tested here explains the Federal Reserve's cap shortfall. Why it does not is not cleanly separable from the population it was tested on. A second qualifier travels with the verdict and is developed where it binds: the multi-vintage stage moved against its own predicted direction, and the inversion is undiagnosed (Appendix~\ref{sec:robustness-extensions}). The paradigm reading this section was written to advance is withdrawn on the strength of that crossing. Section~\ref{sec:robustness-crossdesign} records the withdrawal where the test lives.

\subsection{The Synthetic-Population Estimate: Baseline, Production, Extensions, and the Burnout Falsification}\label{sec:abm-results}

Section~\ref{sec:identification} \emph{declines} that same cross-design recovery as evidence that the imported lock-in elasticity is too small, since a recovery level is not what that design identifies. Read as a scope limit on this section's verdict --- which is how it is used above --- it survives that objection. Neither use licenses the other. The production estimate is that synthetic-population special case, and the case in which the test bites. Every specification tested here moves further from, or fails to close, the \$764.7 billion benchmark (Figure~\ref{fig:recovery}). A purely rational baseline projects a structural cash flow deficit of \$544.0 billion (71.1\% of benchmark). The production specification adds scheduled amortization, a live CUSIP-level SOMA parse, a corrected empirical CPR back-out, a settlement-lag kernel, and the previously excluded 15-year MBS book in cohort weighting and scheduled amortization (structural terms only). Its headline quantity is the seed distribution, not any single draw. The fifty-seed mean is \$103.7 billion, 13.6\% of the benchmark. The frozen seed-42 draw (\$91.0 billion, 11.9\%, near the 34th percentile) is the reproducibility anchor the run manifests freeze. The kernel is retained in production, and both frozen run manifests record it as applied. But it is a documented timing null, not a fit improvement. A pre-fold-in ablation left the peak cross-correlation unchanged at $r = +0.192$ at lag $-3$ with the kernel on or off, and the kernel-applied production run peaks at $+0.195$ at that lag. The production specification is not kernel-free, and every figure reported here is the kernel-applied run's. Seed dispersion dominates the third digit. The seed-level standard deviation is \$24.5 billion ($\pm$1~SD share band 10.4--16.8\%). The 95\% CI of the mean, \$96.9--\$110.5 billion, prices only simulation noise in the mean itself. The 34th percentile is pinned only to within roughly $\pm$7 percentage points from 50 draws. The ABM's share is therefore better read as on the order of 10--17\% than as a three-digit point estimate (Figure~\ref{fig:mc}). The path fit is worse than the level fit and no extension repairs it (Figure~\ref{fig:abmpaths}). Simulated mean CPR is 11.68\% against an empirical 5.14\% (ABM basis). The raw path correlation is negative at lag 0 ($r = -0.318$), and nothing in a $\pm$3-month window exceeds the +0.19 at the three-month lead. The path $R^2$ is $-6.984$. Throughout, $R^2 = 1-\mathrm{SSE}/\mathrm{SST}$ against a mean-only null, so a negative value is a fit worse than a constant at the empirical mean. The claimed object is the aggregate level. The path statistics are reported as a disclosed failure, not as a fitted quantity. No estimator in this paper reproduces the monthly path once the window's shared trend is removed (Section~\ref{sec:pathb}), and no level result anywhere in the paper is scored on path fit. Three literature-grounded behavioral extensions move the model to 54.9\%, then 45.1\%, then 33.7\% of the benchmark (Figure~\ref{fig:waterfall}). They are a debt-to-income constraint with asymmetric loss aversion and a rate-shock freeze, income-scaled curtailment, and multi-vintage coupon composition. The first is a recalibration artifact, not evidence for loss aversion: a real-world turnover floor held fixed under a steeper loss-aversion penalty mechanically inflates rate sensitivity. None touches the CPR path's central failure. Each acts on the level of simulated mobility, not its timing. A fourth variant fails more informatively. Path-dependent prepayment ``burnout,'' implemented as survivor selection in a closed population of 10,000 agents, collapsed simulated CPR to approximately zero and raised trapped liquidity to 147\% of the benchmark, because that population depletes its mobile mass in the low-rate 2020--2021 window before QT begins. That is an absorbing state incompatible with a real MBS pool that continuously receives new originations, and \citepos{lesniewski2026} selection identity in its degenerate limit. It is a negative finding, not used in production, and why burnout enters Section~\ref{sec:hazard} at the cohort level. Appendices~\ref{sec:method-abm}, \ref{sec:method-frictions}, \ref{app:abmspec}, \ref{sec:robustness-extensions} and~\ref{sec:robustness-pipeline} specify the decision rule, every behavioral parameter and its source, the macroeconomic-friction and structural extensions, the three behavioral stages' mechanisms, and the 15-year fold-in. The replication package (its run ledger) reproduces them end to end. This section reports results only.

\begin{figure}[H]
\centering
\includegraphics[width=\textwidth]{fig1_recovery_by_estimator}
\caption{Benchmark recovery by estimator, on the shared accounting basis (standalone-scorer values in parentheses, both bases collected in Table~\ref{tab:bases}; ABM-family rows run through that layer natively). The production ABM row shows the fifty-seed mean (open diamond) with its $\pm 1$~SD band alongside the frozen seed-42 draw (solid dot); Path A's dot is an aggregate-context, non-corroborating exhibit excluded from headline figures (Section~\ref{sec:patha}). Raw lag-0 correlations are reference-only (figure footer; Section~\ref{sec:pathb}); level recovery alone does not rank models.}
\label{fig:recovery}
\end{figure}

\begin{figure}[H]
\centering
\includegraphics[width=\textwidth]{fig3_abm_waterfall}
\caption{The ABM stress-test bridge. Each stage applies the next correction or extension to the household-choice model, and every pre-committed behavioral extension moves the recovery away from, not toward, the benchmark (the multi-vintage stage against its own expected direction, Appendix~\ref{sec:robustness-extensions}): rational baseline 71.1\% $\to$ behavioral gates 54.9\% (a recalibration artifact, Appendix~\ref{sec:robustness-extensions}) $\to$ income-scaled curtailment 45.1\% $\to$ multi-vintage composition 33.7\% $\to$ production corrections 13.2\% (a data-and-accounting stage, not a behavioral one) $\to$ 15-year fold-in 11.9\% (the production headline of Section~\ref{sec:abm-results}; the fold-in itself is worth $-1.3$ points) $\to$ native 15-year gate 11.1\% (the Berger-run calibration-stage variant of Section~\ref{sec:abm-danish}). Stage levels from the frozen run manifests catalogued in the replication package's run ledger.}
\label{fig:waterfall}
\end{figure}

\subsection{Interpretation}\label{sec:abm-interp}

Every ABM specification tested here widens the gap to the benchmark or leaves it unchanged, and its CPR path stays uncorrelated or negatively correlated with the empirical path throughout. Section~\ref{sec:robustness}'s scale-sensitivity replication finds no size effect on the production specification. Two possibilities remain: either household-level decision-making does not govern the majority of the shortfall, or this specification is missing a mechanism. The cross-design leg above bears directly on that either-or without settling it. Real structural covariates move the same rule's recovery from 13.6\% to 60.2\% averaged over fifty seeds. That looks like a missing-population problem rather than a missing-mechanism one. But the behavioral layer stays synthetic in both legs, so the alternative is narrowed, not eliminated. The clean falsification bypasses that free parameter entirely. The simulated-minimum-distance two-moment fit (Section~\ref{sec:robustness-crossdesign}) returns joint infeasibility under a verdict vocabulary fixed before execution. Its nearest admissible fit overshoots the benchmark at 107.1\%, and it does not run through the recalibration scale that qualifies the cross-design reading. Section~\ref{sec:hazard} tests the first branch directly, on a different data source and modeling paradigm. Whichever explanation holds, these results weigh against household mobility choices, rational or behavioral, as the primary price the Federal Reserve paid for secondary-market liquidity. The institutional counterpart is the securitization-policy alternative examined in Section~\ref{sec:abm-danish}: whether the par-payoff rule itself, not household behavior, drives the trap. One asymmetry deserves statement. Section~\ref{sec:identification} insists that this design identifies a differential and not a level, yet I compare the ABM with Path B on levels, because the ABM admits no counterpart to the $\beta_1 = 0$ null. The reason is structural. In \eqref{eq:pathB} the elasticity is a separable term over an independently measured floor, and that floor survives the elasticity's removal. The null therefore stays anchored to observed baseline turnover, and lands at 85.7\% of the benchmark. The ABM has no such term. Its 4--5\% turnover floor is not a component of \eqref{eq:abmmove} but a \emph{product} of it, found by searching the mobility-desire scale $\theta$ until the penalty-bearing rule returns 4--5\% CPR at an 8\% market rate. Zeroing the penalty therefore removes the floor along with the channel. With the production $\theta = 43{,}882.8$ held fixed, the resulting leg returns 40.7\% CPR at the calibration point against the production rule's 4.8\%. It over-retires the book and scores $-259.4$\% of the benchmark. Re-deriving $\theta$ under the null restores the anchor to 4.7\% but scores $+116.3$\%, because the re-search fits the null to the very turnover level the differential is meant to measure. The two readings imply marginals of $+270.4$ and $-105.3$ points: they do not bracket a magnitude; they disagree in sign, and neither is interpretable. I therefore report no ABM differential, and the estimator contrast of Section~\ref{sec:hazard} stays a contrast of levels, with the qualification that entails (run \texttt{abm\_null\_feasibility}; a feasibility probe, not a pre-committed estimate, and nothing in this paper's ranges depends on it).